\section{Bernoulli 过程与离散到达}
\label{sec:06-Random-Processes/02-Counting-Processes/01}

你要统计一台服务的访问量，手上没有埋点，只有一个每秒执行一次的探针：它去问一句``刚过去的这一秒有没有新请求''，然后记下 0 或 1。时间被切成一格一格，每格的结果只有两种。

粗糙，但先别急着嫌弃。这种粒度反而把两个问题的关系摆得格外清楚：一个是``前 100 秒来了多少个请求''，另一个是``第 3 个请求出现在第几秒''。它们看上去是两道题，答案却出自同一张表。

记第 \(n\) 格的指示变量为 \(X_n\)，有事件取 1、无事件取 0。若各格互相独立、每格取 1 的概率都是 \(p\)，序列 \(\{X_n\}\) 就叫 Bernoulli 过程。它大概是最简单的随机过程了：一串独立同分布的硬币。

\subsection{一张图上的两种量法}

把探针的记录画出来。

\begin{center}
\begin{tikzpicture}[x=1cm,y=1cm,font=\small,>=stealth]
  \draw[thick,blue!55]
    (0,0.85) -- (2.1,0.85) -- (2.1,1.15) -- (2.8,1.15) -- (2.8,1.45)
    -- (5.6,1.45) -- (5.6,1.75) -- (7.7,1.75) -- (7.7,2.05) -- (8.4,2.05);
  \draw[->,black!55] (0,0.75) -- (0,2.35) node[right,font=\footnotesize] {\(N_n\)};
  \node[right,font=\footnotesize,black!70] at (8.5,2.05) {\(N_{12}=4\)};
  \foreach \i in {0,1,2,3,4,5,6,7,8,9,10,11}
    \draw[black!40] (0.7*\i,0) rectangle (0.7*\i+0.7,0.45);
  \foreach \i in {2,3,7,10}
    \fill[blue!45] (0.7*\i+0.03,0.03) rectangle (0.7*\i+0.67,0.42);
  \node[below,font=\footnotesize,black!60] at (4.2,-0.05) {第 1 秒到第 12 秒};
  \draw[black!60] (0,-0.55) -- (0,-0.72) -- (2.1,-0.72) -- (2.1,-0.55);
  \node[below,font=\footnotesize] at (1.05,-0.70) {\(W_1=3\)};
  \draw[black!60] (2.1,-0.55) -- (2.1,-0.72) -- (2.8,-0.72) -- (2.8,-0.55);
  \node[below,font=\footnotesize] at (2.45,-0.70) {\(W_2=1\)};
  \draw[black!60] (2.8,-0.55) -- (2.8,-0.72) -- (5.6,-0.72) -- (5.6,-0.55);
  \node[below,font=\footnotesize] at (4.2,-0.70) {\(W_3=4\)};
  \draw[black!60] (5.6,-0.55) -- (5.6,-0.72) -- (7.7,-0.72) -- (7.7,-0.55);
  \node[below,font=\footnotesize] at (6.65,-0.70) {\(W_4=3\)};
  \draw[black!60] (0,-1.45) -- (0,-1.62) -- (5.6,-1.62) -- (5.6,-1.45);
  \node[below,font=\footnotesize] at (2.8,-1.60) {\(T_3=W_1+W_2+W_3=8\)};
\end{tikzpicture}

\emph{同一串探针记录，三种读法。中间是原始的 0/1 序列，涂色格表示有事件；上方的阶梯是累计计数 \(N_n\)，每有一个事件就抬一级；下方是相邻事件之间的间隔 \(W_r\)，把前 \(r\) 段间隔加起来就是第 \(r\) 个事件的到达时刻 \(T_r\)。往后连续时间的图和这张几乎一模一样，只是格子宽度趋于零。}
\end{center}

阶梯与间隔是同一条信息的两种编码，谁也不比谁多。这一点在离散情形下几乎是显然的，但正是它撑起了后面所有的结果，所以值得先在这里坐实。

\subsection{固定时间：计数是 Binomial}

累计计数

\[
N_n=\sum_{k=1}^{n}X_k
\]

是 \(n\) 次独立试验的成功次数，所以 \(N_n\sim\Binomial(n,p)\)，均值 \(np\)、方差 \(np(1-p)\)。均值随时间线性增长，符合直觉。留意方差：\(p\) 很小的时候 \(1-p\approx1\)，方差几乎等于均值。这个巧合不是巧合，下一节会看到它是 Poisson 的标志。

区间 \((m,n]\) 内的增量

\[
N_n-N_m=\sum_{k=m+1}^{n}X_k
\]

只用到落在该区间内的那些格子。于是两件事同时成立：增量的分布只取决于区间长度 \(n-m\)，与起点无关（平稳增量）；不相交区间用的是不相交的格子，因而互相独立（独立增量）。两条性质在这里都是从``每格独立同分布''直接继承下来的，看起来平凡。它们值得单独命名，是因为下一节把格子宽度压到零之后，能活下来的恰好只有这两条加上一条关于瞬时的约束。

\subsection{固定事件数：等待是 Geometric}

现在换个方向问。第 \(r\) 个事件出现在

\[
T_r=\min\{n:N_n=r\}.
\]

这两个方向之间有一座桥，而且是恒等式而非近似：第 \(r\) 个事件在第 \(n\) 格之前（含）出现，当且仅当前 \(n\) 格里至少积够了 \(r\) 个事件。写成事件的相等，

\[
\{T_r\le n\}=\{N_n\ge r\},
\qquad\text{于是}\qquad
P(T_r\le n)=P(N_n\ge r).
\]

左边是等待时间的分布函数，右边是计数分布的尾部和。你只要算清一边，另一边就白送。这座桥不依赖 Bernoulli，任何单调上跳的计数过程都成立，后面还会用到几次。

第一个事件的等待时间容易直接写：前 \(k-1\) 格空着，第 \(k\) 格中了，

\[
P(T_1=k)=(1-p)^{k-1}p,\qquad k=1,2,\ldots
\]

这是 Geometric\((p)\)。推广到第 \(r\) 个事件也不难：前 \(k-1\) 格里恰好放 \(r-1\) 个事件，第 \(k\) 格放第 \(r\) 个，

\[
P(T_r=k)=\binom{k-1}{r-1}p^{r}(1-p)^{k-r},
\]

即负二项分布（negative binomial）。这些计算的细节不重要，重要的是它们全都来自同一句话：格子独立，所以概率直接连乘。

一个事件发生之后，剩下的格子还是原来那串独立试验，什么都没变。所以相邻间隔

\[
W_1,W_2,\ldots\ \overset{\iid}{\sim}\ \mathrm{Geometric}(p),
\qquad T_r=W_1+\cdots+W_r.
\]

反过来也成立：如果一个离散计数过程的间隔是独立同分布的 Geometric\((p)\)，那么每格发生事件的概率必然是 \(p\)，过程只能是 Bernoulli 过程。间隔和过程互相决定，这是本章反复出现的那种对偶的第一次露面。

\subsection{无记忆性在离散情形下的样子}

Geometric 满足

\[
P(W>s+t\given W>s)=P(W>t).
\]

验证只需一行：\(P(W>m)=(1-p)^m\)，代进去左边就是 \((1-p)^{s+t}/(1-p)^{s}\)。含义是，已经空等了 \(s\) 格，剩下还要等几格的分布跟刚开始等的时候完全相同。等待不积累``该轮到我了''的势能。

这条性质是模型能算的根本原因——你不需要记录任何一个事件已经等了多久，当前时刻的状态就够了。它也是最容易被现实推翻的一条。真实系统里等待越久往往越不祥：连着几十秒没有请求，通常不是运气问题，而是上游挂了。第三篇会把这件事展开，并给出一个能在数据上跑的检验。

\subsection{固定的 \texorpdfstring{\(p\)}{p} 撑不了多久}

凌晨三点和晚上八点的访问概率能差一个数量级，把它们混进同一个 \(p\) 里，模型对哪一段都不准。除了速率漂移，还有两类常见的失真：一次点击触发一串后续请求，相邻格子于是相关；限流器攒批放行，事件成堆出现而不是均匀铺开。

失效的方式决定了往哪走。均值随时间漂移，就让 \(p\) 变成 \(p_n\)，对应连续时间里的非齐次过程；同长度窗口的方差远超均值，多半是聚集或异质性，负二项与混合模型是常规解；如果过去的事件会抬高未来的强度，那已经跳出了独立增量的框架，需要自激过程。诊断的价值在于，偏离的方向本身就指出了替代品。

顺带一提，这些偏离在中心区域常常看不出来——平均每分钟多少个请求，谁拟合都差不多。差别出在尾部，而容量规划、超时阈值、SLA 承诺关心的恰好是尾部。

最后一步把格子宽度压到零。若每个小区间发生一次事件的概率与区间长度成正比，而发生两次以上的概率是更高阶的小量，离散逼近就收敛到一个连续时间的过程。下一节从这个极限出发，把它需要的假设减到三条。
